\chapter{The Cooper Pair Ledger}

A Cooper pair is not two separately tagged electrons walking side by side in the ledger. In a superconductor, the
pair belongs to the condensate. Its useful public facts are charge \(2e\), even parity, phase coherence, and a
gap separating the paired ground from quasiparticle excitations. The apparatus therefore records the pair not as
two independent electron rows, but as a boundary state with its own charge step and its own phase cost.

This is where the earlier electron naming has to be handled with care. The electron was named by counting: box
one, charge read as a single negative unit. The Cooper-pair boundary now receives that electron through a second
electron position at the junction, but it still does not erase the distinction between one and two. A pair
transfer moves charge in units of \(2e\). An electron-electron contact carries its own residual. The interface is
exactly the place where those units meet and refuse to become one another.

\section{The Gauge Facts}

The existing Lean gauge already proves the small facts needed at the door. In the companion file they are exposed
at the face level:

\begin{quote}\small
\begin{verbatim}
theorem singleElectron_has_no_tange :
    Face.positronCount .singleElectron = 0 := by
  decide

theorem junctionElectron_has_no_tange :
    Face.positronCount .junctionElectron = 0 := by
  decide

theorem cooperPair_has_one_tange :
    Face.positronCount .cooperPair = 1 := by
  decide

def Face.totalCount (face : Face) : Nat :=
  face.matterCount + face.positronCount

theorem cooperPair_tallies_the_gauge :
    Face.totalCount .cooperPair = Gate.all.length := by
  decide
\end{verbatim}
\end{quote}

These are finite gauge theorems, not microscopic superconductivity. They say how the apparatus reads the faces
once the face has been supplied. The outer electron and the junction electron both sit at the all-funge corner.
The Cooper-pair face sits one tange away. The matter count and the tange count partition the whole gauge; no
hidden fourth class is allowed to drift between them.

This is enough for an instruction manual because the pair computation is not trying to derive the condensate
from first principles. It is trying to add a contact law to the apparatus. A contact law says: when the outer
electron meets the junction electron, carry the electron-electron residual; when the interface is read as a pair,
count the pair transfer and charge the phase; when the run unwinds, discharge those residuals in the reverse
order.

\section{The Physical Scales}

A real device needs physical scales. The companion file leaves a record for them:

\begin{quote}\small
\begin{verbatim}
structure JunctionScales where
  gap : Float
  josephsonEnergy : Float
  chargingEnergy : Float
  gateOffset : Float
deriving Repr
\end{verbatim}
\end{quote}

The gap is the energy cost to break a pair into quasiparticles. The Josephson energy is the strength of coherent
pair transfer across the junction. The charging energy is the price of changing the island charge. The gate
offset records where the hand has biased the island. These four numbers are the first practical inputs for a
Cooper-pair device. They should be entered with units and uncertainties in the laboratory ledger, even though the
minimal Lean skeleton stores only floating-point demonstration slots.

The apparatus grade is mixed again. The gauge facts above are PROVED inside the finite Lean gauge. The physical
scales are GIVEN or MEASURED depending on how the laboratory obtains them. The identification of a particular
boundary event as Josephson tunneling, Andreev reflection, poisoning, or pair breaking is a READING laid over the
device. The computation is honest only when those grades remain separated.

\section{The State Record}

The running state contains only what the interface needs for its first manual:

\begin{quote}\small
\begin{verbatim}
structure InteractionState where
  stage : Stage
  face : Face
  electronElectronEvents : Nat
  electronElectronResidual : Int
  pairTransfers : Int
  singleEvents : Nat
  pairEvents : Nat
  phaseTurns : Int
  costSpent : Nat
deriving Repr
\end{verbatim}
\end{quote}

The \(\texttt{stage}\) says where the four-stage trace has stopped. The \(\texttt{face}\) says how that stage
reads through the finite gauge. The two electron-electron fields count the outer/junction contacts and carry a
signed residual until the return step unwinds it. The \(\texttt{pairTransfers}\) field is signed because the
interface can move a pair into or out of the condensate. The \(\texttt{singleEvents}\) and
\(\texttt{pairEvents}\) fields are unsigned counts of how many electron and pair faces have been seen. The
\(\texttt{phaseTurns}\) field is signed because a pair transfer can advance or unwind the simplified phase
ledger. The \(\texttt{costSpent}\) field is the operator's clock.

The initial state starts at the outer electron:

\begin{quote}\small
\begin{verbatim}
def initial : InteractionState :=
  { stage := .outerElectron
    face := .singleElectron
    electronElectronEvents := 0
    electronElectronResidual := 0
    pairTransfers := 0
    singleEvents := 1
    pairEvents := 0
    phaseTurns := 0
    costSpent := 0 }
\end{verbatim}
\end{quote}

Starting at the outer electron is a choice, not a law. A different experiment may start at the junction electron,
at the pair face, or at a state already carrying phase. The manual starts here because the user asked how to put
in the electron \(\to\) electron \(\to\) Cooper-pair interaction. The outer electron is brought to the door
first. The junction electron answers next. The condensate answers after that.
